\documentclass[a4paper,10pt,brazil]{article}
\usepackage{provastyle}
\usepackage{menukeys}

\begin{document}
\makeexamtitlepage{UFOP}{BCC222: Programação Funcional}{José Romildo}{2026/1}{Primeira Prova}{05}
\begin{multicols}{2}
\questionTitle{1}{Acesso composto a pares}
\begin{flushright}
  \footnotesize
  \textit{\textbf{6: Estruturas de dados básicas}}\\
  \textit{c06-tuplas-acesso-01}\\

\end{flushright}
Considere a expressão abaixo.
\begin{minted}{haskell}
snd (fst (('x', 19), "ok"))
\end{minted}
Qual é o seu resultado?

\begin{enumerate}
  \item \haskellinline{ ('x',19) }

  \item A expressão resulta em erro de compilação.

  \item \haskellinline{ 'x' }

  \item \haskellinline{ 19 }

  \item \haskellinline{ "ok" }


\end{enumerate}

\questionTitle{2}{A sintaxe de múltiplas linhas}
\begin{flushright}
  \footnotesize
  \textit{\textbf{2: Ambiente de desenvolvimento}}\\
  \textit{c02-fluxo-trabalho-03}\\

\end{flushright}
Ao tentar copiar uma definição funcional de 4 linhas de um PDF e colar diretamente no \textbf{prompt} do GHCi, o terminal costuma gerar erros acusando que identificadores não estão em escopo, pois avalia cada linha independentemente. 

Qual é a solução mais segura, utilizando delimitadores, para que o GHCi aceite todo o bloco colado como uma única construção sintática?

\begin{enumerate}
  \item Utilizar aspas triplas \texttt{"""} tanto no topo quanto na base do trecho copiado.

  \item Adicionar barras duplas \haskellinline{//} no fim de cada linha para avisar o interpretador de que o comando continua na próxima.

  \item Inserir o comando de avaliação em lote \texttt{:batch} na linha imediatamente anterior.

  \item Involucrar o bloco copiado digitando \texttt{:\{} antes de colar, e \texttt{:\}} ao finalizar.

  \item Envolver a expressão utilizando parênteses duplos \texttt{((} no início e \texttt{))} no fim.


\end{enumerate}

\questionTitle{3}{Defaulting com show e literal fracionário}
\begin{flushright}
  \footnotesize
  \textit{\textbf{7: Polimorfismo}}\\
  \textit{cap07-defaulting-show-fracionario}\\

\end{flushright}
Em um arquivo Haskell sem \haskellinline{ExtendedDefaultRules}, considere:

\begin{minted}{haskell}
u = show 4.12
\end{minted}

Por que essa definição pode ser aceita sem assinatura explícita?

\begin{enumerate}
  \item Porque \haskellinline{show} sempre converte literais fracionários para \haskellinline{Float}.
  \item Porque toda ambiguidade com \haskellinline{read} e \haskellinline{show} é resolvida automaticamente pelo Relatório Haskell.
  \item Porque \haskellinline{Show} é uma classe numérica e força o tipo \haskellinline{Integer}.
  \item Porque \haskellinline{4.12} tem tipo fixo \haskellinline{String}.
  \item A variável de tipo ocorre apenas nas restrições \haskellinline{Fractional} e \haskellinline{Show}; como há uma restrição numérica, a lista padrão escolhe \haskellinline{Double}.

\end{enumerate}

\questionTitle{4}{Refatoração para forma mais concisa}
\begin{flushright}
  \footnotesize
  \textit{\textbf{5: Expressão condicional}}\\
  \textit{c05-if-vs-guardas-02}\\

\end{flushright}
Considere a função:
\begin{haskellcode}
classifica n
  | n > 0     = "Positivo"
  | n < 0     = "Negativo"
  | n == 0    = "Zero"
  | otherwise = "Impossível"
\end{haskellcode}
Qual definição é funcionalmente equivalente, mas mais concisa?

\begin{enumerate}
  \item \begin{haskellcode}
classifica n
  | n == 0    = "Zero"
  | otherwise = "Positivo ou Negativo"
\end{haskellcode}

  \item \begin{haskellcode}
classifica n
  | n > 0     = "Positivo"
  | otherwise = "Negativo"
\end{haskellcode}

  \item \begin{haskellcode}
classifica n
  | n >= 0    = "Positivo"
  | otherwise = "Negativo"
\end{haskellcode}

  \item \begin{haskellcode}
classifica n
  | n > 0     = "Positivo"
  | n < 0     = "Negativo"
  | otherwise = "Zero"
\end{haskellcode}

  \item A função original não pode ser simplificada sem mudar seu comportamento.


\end{enumerate}

\questionTitle{5}{A primeira linguagem funcional prática}
\begin{flushright}
  \footnotesize
  \textit{\textbf{1: Paradigmas de programação}}\\
  \textit{c01-historia-evolucao-02}\\

\end{flushright}
Criada na virada da década de 1950 por John McCarthy no MIT, qual linguagem de programação é historicamente consagrada como o primeiro esforço prático para incorporar o paradigma funcional, introduzindo conceitos inovadores à época, como a recursão nativa e o \emph{Garbage Collection}?

\begin{enumerate}
  \item Lisp

  \item Fortran

  \item ML (MetaLanguage)

  \item Miranda

  \item COBOL


\end{enumerate}

\questionTitle{6}{Assinatura mais geral e uso excessivo de restrições}
\begin{flushright}
  \footnotesize
  \textit{\textbf{7: Polimorfismo}}\\
  \textit{cap07-param-funcao-restrita-vs-geral}\\

\end{flushright}
Um estudante propôs a seguinte assinatura para a função abaixo:

\begin{minted}{haskell}
duplicaPar x = (x, x)

duplicaPar :: Eq a => a -> (a, a)
\end{minted}

Qual é a melhor avaliação técnica?

\begin{enumerate}
  \item A restrição \haskellinline{Eq a} é obrigatória, pois repetir um valor exige compará-lo consigo mesmo.
  \item A restrição \haskellinline{Eq a} é desnecessária; a assinatura mais geral é \mintinline{haskell}{duplicaPar :: a -> (a, a)}.
  \item A assinatura correta é \mintinline{haskell}{duplicaPar :: a -> a}, pois a tupla não altera o tipo do valor.
  \item A função é monomórfica, pois o mesmo \haskellinline{x} aparece duas vezes no corpo.
  \item A função deveria exigir \haskellinline{Ord a}, pois tuplas sempre precisam ser ordenáveis.

\end{enumerate}

\questionTitle{7}{Equivalência com True}
\begin{flushright}
  \footnotesize
  \textit{\textbf{5: Expressão condicional}}\\
  \textit{c05-otherwise-02}\\

\end{flushright}
Qual definição é equivalente ao uso usual de \haskellinline|otherwise| na última guarda?

\begin{enumerate}
  \item Substituir \haskellinline|otherwise| por \haskellinline|else|.

  \item Substituir \haskellinline|otherwise| por \haskellinline|False|.

  \item Substituir \haskellinline|otherwise| por \haskellinline|True|.

  \item Substituir \haskellinline|otherwise| por \haskellinline|undefined|.

  \item Substituir \haskellinline|otherwise| por \haskellinline|()|.


\end{enumerate}

\questionTitle{8}{Comando para recarregar um arquivo no GHCi}
\begin{flushright}
  \footnotesize
  \textit{\textbf{3: Expressões e definições}}\\
  \textit{c03-ghci-definicoes-03}\\

\end{flushright}
Após carregar um arquivo Haskell no GHCi e editar o código no editor, qual comando recarrega o último arquivo carregado?

\begin{enumerate}
  \item \haskellinline|:edit|, pois ele salva e recompila automaticamente o módulo.

  \item \haskellinline|:type|, pois ele recompila e mostra os tipos de todas as definições.

  \item \haskellinline|:reload|, ou sua abreviação \haskellinline|:r|.

  \item \haskellinline|:load| sem argumento, pois ele sempre recarrega o arquivo anterior.

  \item \haskellinline|:browse|, pois ele atualiza o arquivo fonte em disco.


\end{enumerate}

\questionTitle{9}{Valor produzido por uma ação não é impresso automaticamente}
\begin{flushright}
  \footnotesize
  \textit{\textbf{10: Ações de E/S recursivas}}\\
  \textit{c10-return-tipos-04}\\

\end{flushright}
Analise o programa:

\begin{haskellcode}
main :: IO ()
main = do
  return 42
  putStrLn "fim"
\end{haskellcode}

O que ocorre quando \haskellinline|main| é executado?

\begin{enumerate}
  \item O número \haskellinline|42| é impresso e, em seguida, a palavra \haskellinline|"fim"| é impressa.

  \item A execução termina em \haskellinline|return 42| e a linha \haskellinline|putStrLn "fim"| nunca é executada.

  \item O programa não compila, pois uma ação \haskellinline|IO Int| não pode aparecer antes de uma ação \haskellinline|IO ()| em um bloco \haskellinline|do|.

  \item A ação \haskellinline|return 42| é executada e seu resultado é descartado; depois, \haskellinline|"fim"| é impresso.

  \item O valor \haskellinline|42| é armazenado automaticamente em uma variável chamada \haskellinline|it| dentro de \haskellinline|main|.


\end{enumerate}

\questionTitle{10}{Inferência com ceiling e divisão fracionária}
\begin{flushright}
  \footnotesize
  \textit{\textbf{7: Polimorfismo}}\\
  \textit{cap07-inferencia-num-degraus}\\

\end{flushright}
Considere a função:

\begin{minted}{haskell}
numDegraus alturaDegrau alturaDesejada =
  ceiling (alturaDesejada / alturaDegrau)
\end{minted}

Qual é a assinatura de tipo mais geral?

\begin{enumerate}
  \item \mintinline{haskell}{numDegraus :: Double -> Double -> Double}
  \item \mintinline{haskell}{numDegraus :: (Num a, Ord a) => a -> a -> Bool}
  \item \mintinline{haskell}{numDegraus :: (RealFrac a, Integral b) => a -> a -> b}
  \item \mintinline{haskell}{numDegraus :: Fractional a => a -> a -> a}
  \item \mintinline{haskell}{numDegraus :: Integral a => a -> a -> a}

\end{enumerate}

\questionTitle{11}{Escolha entre arredondar e converter integralmente}
\begin{flushright}
  \footnotesize
  \textit{\textbf{7: Polimorfismo}}\\
  \textit{cap07-conversao-arredondamento-escolha}\\

\end{flushright}
Uma função deve receber um \haskellinline{Double} e produzir um \haskellinline{Int} arredondando para o inteiro mais próximo. Qual função é adequada?
\begin{enumerate}
  \item \haskellinline{fromInteger}
  \item \haskellinline{round}
  \item \haskellinline{fromIntegral}
  \item \haskellinline{realToFrac}, pois sempre retorna um tipo integral
  \item \haskellinline{show}, pois a conversão passa por texto

\end{enumerate}

\questionTitle{12}{Classe mínima para igualdade e ordenação}
\begin{flushright}
  \footnotesize
  \textit{\textbf{7: Polimorfismo}}\\
  \textit{cap07-adhoc-classe-minima-comparacao}\\

\end{flushright}
Considere a função:

\begin{minted}{haskell}
entre x y z = x <= y && y <= z
\end{minted}

Qual restrição de classe é suficiente e mais geral para o tipo dos argumentos?

\begin{enumerate}
  \item \haskellinline{Show a}, pois o operador \haskellinline{(&&)} precisa imprimir os operandos.
  \item \haskellinline{Num a}, pois qualquer comparação exige que os valores sejam numéricos.
  \item Nenhuma restrição é necessária, pois \haskellinline{(<=)} é paramétrico puro.
  \item \haskellinline{Eq a}, pois igualdade é suficiente para implementar \haskellinline{(<=)}.
  \item \haskellinline{Ord a}, pois \haskellinline{(<=)} pertence a \haskellinline{Ord} e \haskellinline{Ord} já pressupõe \haskellinline{Eq}.

\end{enumerate}

\questionTitle{13}{Valores e forma normal}
\begin{flushright}
  \footnotesize
  \textit{\textbf{3: Expressões e definições}}\\
  \textit{c03-comentarios-valores-03}\\

\end{flushright}
No contexto da avaliação de expressões em Haskell, qual alternativa caracteriza corretamente a relação entre expressões e valores?

\begin{enumerate}
  \item Uma expressão só é considerada valor se tiver sido associada a um nome por uma definição.

  \item Toda expressão escrita pelo programador já é imediatamente um valor, mesmo antes de qualquer avaliação.

  \item Valores são apenas números; listas, tuplas e funções não são considerados valores em Haskell.

  \item Todo valor é uma expressão, mas nem toda expressão já é um valor; uma expressão torna-se um valor quando atinge uma forma que não pode mais ser reduzida.

  \item Uma expressão em forma normal ainda precisa ser reduzida pelo menos uma vez para produzir um valor.


\end{enumerate}

\questionTitle{14}{Ordem das definições locais em where}
\begin{flushright}
  \footnotesize
  \textit{\textbf{3: Expressões e definições}}\\
  \textit{c03-definicoes-where-03}\\

\end{flushright}
Considere a definição:

\begin{haskellcode}
g x = a * b
  where
    b = x + 2
    a = b - 1
\end{haskellcode}

Qual é o valor de \haskellinline|g 4|?

\begin{enumerate}
  \item Ocorre erro, pois definições em \haskellinline|where| não podem depender umas das outras.

  \item \haskellinline|20|.

  \item \haskellinline|30|.

  \item Ocorre erro, pois \haskellinline|a| aparece antes de \haskellinline|b| no corpo da função.

  \item \haskellinline|18|.


\end{enumerate}

\questionTitle{15}{Análise de assinatura de tipo}
\begin{flushright}
  \footnotesize
  \textit{\textbf{4: Tipos de Dados}}\\
  \textit{c04-assinatura-analise-01}\\

\end{flushright}
Qual das seguintes definições de função corresponde à assinatura de tipo \haskellinline{String -> Int -> Bool}?

\begin{enumerate}
  \item \begin{haskellcode}
verifica s = s == "OK"
\end{haskellcode}

  \item \begin{haskellcode}
verifica s n = s ++ show n
\end{haskellcode}

  \item \begin{haskellcode}
verifica s n = length s > n
\end{haskellcode}

  \item \begin{haskellcode}
verifica b s = not b && null s
\end{haskellcode}

  \item \begin{haskellcode}
verifica n s = n > 0 && s == "OK"
\end{haskellcode}


\end{enumerate}

\questionTitle{16}{Boas práticas em assinaturas explícitas}
\begin{flushright}
  \footnotesize
  \textit{\textbf{7: Polimorfismo}}\\
  \textit{cap07-comparacao-assinaturas-explicitas}\\

\end{flushright}
Mesmo com inferência de tipos, por que é boa prática escrever assinaturas explícitas para funções de topo?
\begin{enumerate}
  \item Porque Haskell não consegue inferir tipos de funções definidas pelo programador.
  \item Porque a assinatura documenta a intenção, ajuda a localizar erros e pode impedir que uma função fique mais geral ou mais restrita do que o planejado.
  \item Porque assinaturas explícitas desativam a checagem de tipos, facilitando a compilação.
  \item Porque assinaturas explícitas tornam as funções mutáveis.
  \item Porque toda função sem assinatura passa a ter tipo \haskellinline{String}.

\end{enumerate}

\questionTitle{17}{Tipo de um bloco do}
\begin{flushright}
  \footnotesize
  \textit{\textbf{8: Programas interativos}}\\
  \textit{c08-tipos-acoes-04}\\

\end{flushright}
Qual é o tipo da ação abaixo?

\begin{minted}{haskell}
acao = do
  putStrLn "Digite uma palavra:"
  getLine
\end{minted}

\begin{enumerate}
  \item \haskellinline|String|, pois o bloco \haskellinline|do| extrai automaticamente o valor produzido por \haskellinline|getLine|.

  \item \haskellinline|()|, pois o valor lido por \haskellinline|getLine| é descartado.

  \item \haskellinline|IO ()|, pois o bloco contém uma chamada a \haskellinline|putStrLn|.

  \item \haskellinline|IO String|, pois a última ação do bloco é \haskellinline|getLine|.

  \item \haskellinline|String -> IO String|, pois a primeira linha imprime uma mensagem.


\end{enumerate}

\questionTitle{18}{Comparação entre soma tardia e soma acumulada}
\begin{flushright}
  \footnotesize
  \textit{\textbf{10: Ações de E/S recursivas}}\\
  \textit{c10-acumuladores-cauda-04}\\

\end{flushright}
Compare as duas definições:

\begin{haskellcode}
-- Versão A
lerA :: IO Integer
lerA = do
  n <- readLn
  if n == 0
    then return 0
    else do
      s <- lerA
      return (n + s)

-- Versão B
lerB :: Integer -> IO Integer
lerB total = do
  n <- readLn
  if n == 0
    then return total
    else lerB (total + n)
\end{haskellcode}

Qual alternativa expressa melhor a diferença estrutural entre elas?

\begin{enumerate}
  \item A versão A deixa uma soma pendente após a chamada recursiva; a versão B incorpora o valor lido ao acumulador antes da próxima chamada.

  \item A versão A lê todos os valores antes de testar o sentinela; a versão B testa o sentinela antes de ler.

  \item A versão B ignora os valores lidos, pois não há variável ligada por \haskellinline|<-| no ramo recursivo.

  \item As duas versões têm exatamente a mesma estrutura de chamadas, diferindo apenas nos nomes das variáveis.

  \item A versão A é pura e a versão B é impura, pois apenas B recebe parâmetro.


\end{enumerate}

\questionTitle{19}{Validação após leitura numérica}
\begin{flushright}
  \footnotesize
  \textit{\textbf{8: Programas interativos}}\\
  \textit{c08-conversao-validacao-04}\\

\end{flushright}
Considere a ação:

\begin{minted}{haskell}
main :: IO ()
main = do
  putStrLn "Digite uma idade:"
  idade <- readLn :: IO Int
  if idade < 0
    then putStrLn "Idade inválida!"
    else print idade
\end{minted}

Qual afirmação é correta?

\begin{enumerate}
  \item O programa valida idades negativas depois de converter a entrada para \haskellinline|Int|. Entretanto, uma entrada textual que não represente um inteiro ainda pode causar exceção de conversão.

  \item O programa não compila, pois \haskellinline|readLn| não pode receber anotação de tipo.

  \item O ramo \haskellinline|else| transforma automaticamente a idade em \haskellinline|String|.

  \item O programa valida qualquer entrada inválida, incluindo textos como \haskellinline|"abc"|, sem risco de exceção.

  \item O teste \haskellinline|idade < 0| é executado antes da leitura da entrada.


\end{enumerate}

\questionTitle{20}{Identificadores alfanuméricos válidos}
\begin{flushright}
  \footnotesize
  \textit{\textbf{3: Expressões e definições}}\\
  \textit{c03-arquivos-identificadores-03}\\

\end{flushright}
De acordo com as regras para identificadores alfanuméricos usados como nomes de variáveis ou funções, qual alternativa apresenta um identificador válido?

\begin{enumerate}
  \item \haskellinline|module|.

  \item \haskellinline|ValorFinal|.

  \item \haskellinline|valor-final|.

  \item \haskellinline|2valor|.

  \item \haskellinline|valor_final'|.


\end{enumerate}

\questionTitle{21}{Interpretação de Num a => a}
\begin{flushright}
  \footnotesize
  \textit{\textbf{4: Tipos de Dados}}\\
  \textit{c04-ghci-interpretacao-01}\\

\end{flushright}
Ao executar \texttt{:type 42} no GHCi, a resposta é \haskellinline{42 :: Num a => a}. O que isso significa?

\begin{enumerate}
  \item \haskellinline{42} é uma constante numérica que só pode ser do tipo \haskellinline{Int}.

  \item \haskellinline{42} pode ser interpretado como um valor de qualquer tipo que pertença à classe \haskellinline{Num} (como \haskellinline{Int}, \haskellinline{Integer}, \haskellinline{Double}).

  \item Indica um erro de tipo, pois o tipo é ambíguo e não pode ser resolvido.

  \item \haskellinline{42} é um valor de um tipo desconhecido que não pode ser usado em operações aritméticas.

  \item Significa que \haskellinline{42} é uma função que recebe um \haskellinline{Num} e retorna um \haskellinline{a}.


\end{enumerate}

\questionTitle{22}{Descarte de valores em do}
\begin{flushright}
  \footnotesize
  \textit{\textbf{8: Programas interativos}}\\
  \textit{c08-blocos-do-04}\\

\end{flushright}
No bloco abaixo, a primeira chamada a \haskellinline|getLine| está isolada em uma linha:

\begin{minted}{haskell}
main :: IO ()
main = do
  getLine
  nome <- getLine
  putStrLn nome
\end{minted}

Qual é o comportamento correto?

\begin{enumerate}
  \item O programa imprime as duas linhas lidas, uma em cada linha.

  \item A primeira chamada a \haskellinline|getLine| não é executada, pois seu resultado não é usado.

  \item O programa não compila, pois toda ação em um bloco \haskellinline|do| deve usar \haskellinline|<-|.

  \item A primeira ação \haskellinline|getLine| é executada e consome uma linha da entrada, mas seu resultado é descartado. A segunda leitura é associada a \haskellinline|nome| e depois impressa.

  \item As duas chamadas a \haskellinline|getLine| leem a mesma linha, pois Haskell reutiliza automaticamente ações de E/S.


\end{enumerate}

\questionTitle{23}{Combinação de read e show com anotação}
\begin{flushright}
  \footnotesize
  \textit{\textbf{7: Polimorfismo}}\\
  \textit{cap07-expressao-show-read-anotado}\\

\end{flushright}
Qual é o valor da expressão?

\begin{minted}{haskell}
show (read "10" :: Int)
\end{minted}

\begin{enumerate}
  \item \mintinline{haskell}{"Int"}
  \item \mintinline{haskell}{10}
  \item A expressão é ambígua, mesmo com a anotação \haskellinline{:: Int}.
  \item A expressão tem erro de tipo, pois \haskellinline{show} não aceita \haskellinline{Int}.
  \item \mintinline{haskell}{"10"}

\end{enumerate}

\questionTitle{24}{A distinção entre GHC e GHCi}
\begin{flushright}
  \footnotesize
  \textit{\textbf{2: Ambiente de desenvolvimento}}\\
  \textit{c02-ecossistema-ferramentas-03}\\

\end{flushright}
O conjunto de ferramentas do Glasgow Haskell Compiler fornece dois componentes essenciais: o \texttt{ghc} e o \texttt{ghci}. Qual é a distinção primária entre eles durante o desenvolvimento?

\begin{enumerate}
  \item O \texttt{ghc} traduz funções puras matemáticas, enquanto o \texttt{ghci} é reservado apenas para compilar funções que realizam operações de entrada e saída (I/O).

  \item O \texttt{ghc} gerencia a formatação do código no editor de texto, ao passo que o \texttt{ghci} executa os testes de unidade automatizados.

  \item Não há diferença funcional; ambos são atalhos distintos de terminal que iniciam exatamente a mesma ferramenta de \emph{build}.

  \item O \texttt{ghc} é utilizado exclusivamente no sistema operacional Linux, enquanto o \texttt{ghci} é a versão portada para o Windows.

  \item O \texttt{ghc} é o compilador de linha de comando que gera binários executáveis, enquanto o \texttt{ghci} é o ambiente interativo (\emph{REPL}) para avaliação em tempo real.


\end{enumerate}

\questionTitle{25}{Decisão em programa interativo}
\begin{flushright}
  \footnotesize
  \textit{\textbf{8: Programas interativos}}\\
  \textit{c08-programas-completos-04}\\

\end{flushright}
Considere:

\begin{minted}{haskell}
main :: IO ()
main = do
  nota <- readLn :: IO Double
  if nota >= 6.0
    then putStrLn "Aprovado"
    else putStrLn "Reprovado"
\end{minted}

Se o usuário digitar \haskellinline|5.5|, o que será impresso?

\begin{enumerate}
  \item \haskellinline|Aprovado|

  \item \haskellinline|False|

  \item O programa não compila, pois \haskellinline|if| não pode aparecer dentro de \haskellinline|do|.

  \item \haskellinline|5.5|

  \item \haskellinline|Reprovado|


\end{enumerate}

\questionTitle{26}{Inferência de tipo de isDigit}
\begin{flushright}
  \footnotesize
  \textit{\textbf{4: Tipos de Dados}}\\
  \textit{c04-inferencia-basica-02}\\

\end{flushright}
Considere \haskellinline{import Data.Char}. Qual é o tipo inferido para a expressão \haskellinline{isDigit '5'}?

\begin{enumerate}
  \item \haskellinline{Char}

  \item \haskellinline{Bool}

  \item \haskellinline{True}

  \item \haskellinline{Num a => a}

  \item \haskellinline{Char -> Bool}


\end{enumerate}

\questionTitle{27}{Turma vazia no fechamento de notas}
\begin{flushright}
  \footnotesize
  \textit{\textbf{10: Ações de E/S recursivas}}\\
  \textit{c10-bordas-robustez-04}\\

\end{flushright}
Em um programa de fechamento de notas, suponha que os dados dos alunos sejam lidos em uma lista e que o percentual de aprovados seja calculado por uma função pura que divide pelo tamanho da lista. Qual cuidado é necessário?

\begin{enumerate}
  \item É necessário tratar explicitamente a turma vazia, evitando divisão por zero e definindo uma mensagem ou convenção adequada para esse caso.

  \item A função deve chamar \haskellinline|readLn| novamente se a lista estiver vazia, mesmo sendo uma função pura.

  \item Nenhum cuidado é necessário, pois \haskellinline|length []| é automaticamente convertido para \haskellinline|1| em divisões.

  \item A lista vazia deve ser substituída por uma lista infinita, para que a média da turma exista.

  \item O problema só ocorre se as matrículas forem negativas; o tamanho da lista não interfere nos percentuais.


\end{enumerate}

\questionTitle{28}{A ação getContents}
\begin{flushright}
  \footnotesize
  \textit{\textbf{8: Programas interativos}}\\
  \textit{c08-entrada-saida-04}\\

\end{flushright}
Qual alternativa descreve corretamente a ação \haskellinline|getContents|?

\begin{enumerate}
  \item \haskellinline|getContents :: IO ()| é uma ação que descarta toda a entrada padrão.

  \item \haskellinline|getContents :: IO String| é uma ação que lê todo o conteúdo disponível na entrada padrão como uma única string, usualmente de forma preguiçosa.

  \item \haskellinline|getContents :: String| é uma string pura contendo o texto digitado antes da execução do programa.

  \item \haskellinline|getContents :: FilePath -> IO String| exige sempre o nome de um arquivo como argumento.

  \item \haskellinline|getContents :: IO Char| lê apenas o primeiro caractere digitado pelo usuário.


\end{enumerate}

\questionTitle{29}{Análise declarativa da abstração do Quicksort}
\begin{flushright}
  \footnotesize
  \textit{\textbf{1: Paradigmas de programação}}\\
  \textit{c01-mecanismo-reducao-02}\\

\end{flushright}
Considere a cláusula principal da implementação declarativa do \emph{Quicksort} em Haskell:

\begin{haskellcode}
qsort (pivo:resto) = qsort (filter (<= pivo) resto) ++ 
                     [pivo] ++ 
                     qsort (filter (> pivo) resto)
\end{haskellcode}

Ao comparar este trecho com a tradicional implementação iterativa em C, qual das seguintes afirmações avalia corretamente a diferença de abstração utilizada na partição dos elementos?

\begin{enumerate}
  \item A versão funcional abstrai a manipulação de índices de memória e laços \cinline|while|, descrevendo a lista particionada como o resultado da aplicação de funções de filtragem baseadas em propriedades lógicas da lista.

  \item A versão funcional é inerentemente mais ineficiente pois exige que o compilador traduza o operador de concatenação (\haskellinline|++|) em loops duplamente encadeados estritos.

  \item A versão funcional oculta a passagem de ponteiros de memória nos bastidores do operador \haskellinline|:|, realizando uma mutação de array idêntica à troca de variáveis (\cinline|swap|) do algoritmo em C.

  \item A versão funcional falha em implementar a estratégia de \emph{divisão e conquista}, pois a linguagem Haskell não suporta chamadas de recursão de cauda nativa como a linguagem C.

  \item A versão funcional requer que a lista de entrada seja previamente ordenada, caso contrário a função \haskellinline|filter| lançará uma exceção de limite de tempo de execução.


\end{enumerate}

\questionTitle{30}{Classificação da forma de recursão}
\begin{flushright}
  \footnotesize
  \textit{\textbf{9: Funções recursivas}}\\
  \textit{c09-formas-classificacao-01}\\

\end{flushright}
Qual alternativa classifica corretamente as formas de recursão estudadas no capítulo?

\begin{enumerate}
  \item Toda função com mais de uma guarda é, necessariamente, uma função mutuamente recursiva.

  \item \haskellinline|fatorial| é recursão simples; \haskellinline|fib| é recursão múltipla; \haskellinline|par| e \haskellinline|impar| formam recursividade mútua.

  \item \haskellinline|fatorial| e \haskellinline|fib| são recursividade mútua, pois ambas têm casos base.

  \item \haskellinline|fatorial| é recursão múltipla; \haskellinline|fib| é recursão de cauda; \haskellinline|par| e \haskellinline|impar| não são recursivas.

  \item \haskellinline|fib| não é recursiva, pois usa soma no caso recursivo.


\end{enumerate}

\questionTitle{31}{Entrada que provoca falha}
\begin{flushright}
  \footnotesize
  \textit{\textbf{5: Expressão condicional}}\\
  \textit{c05-guardas-exaustividade-02}\\

\end{flushright}
Para qual chamada a função abaixo falha em tempo de execução?
\begin{haskellcode}
comparar x y
  | x > y = "Maior"
  | x < y = "Menor"
\end{haskellcode}

\begin{enumerate}
  \item \haskellinline|comparar 5 10|

  \item \haskellinline|comparar 10 5|

  \item \haskellinline|comparar 0 (-1)|

  \item Nenhuma, pois o compilador adiciona \haskellinline|otherwise| implicitamente.

  \item \haskellinline|comparar 5 5|


\end{enumerate}

\questionTitle{32}{Avaliação com cálculo local}
\begin{flushright}
  \footnotesize
  \textit{\textbf{5: Expressão condicional}}\\
  \textit{c05-where-avaliacao-02}\\

\end{flushright}
Considere a função:
\begin{haskellcode}
analisa y
  | x > 10    = "Alto"
  | otherwise = "Baixo"
  where
    x = y * y + 1
\end{haskellcode}
Qual é o resultado de \haskellinline|analisa 3|?

\begin{enumerate}
  \item Ocorre um erro de tipo em \haskellinline|x > 10|.

  \item \haskellinline|10|

  \item \haskellinline|3|

  \item \haskellinline|"Alto"|

  \item \haskellinline|"Baixo"|


\end{enumerate}

\questionTitle{33}{A Crise do Software e a solução funcional}
\begin{flushright}
  \footnotesize
  \textit{\textbf{1: Paradigmas de programação}}\\
  \textit{c01-engenharia-software-02}\\

\end{flushright}
O termo \emph{Crise do Software} surgiu historicamente devido à dificuldade de entregar sistemas complexos no prazo, dentro do orçamento e isentos de falhas catastróficas em tempo de execução. Como o paradigma funcional atua para mitigar diretamente a raiz estrutural dessa crise?

\begin{enumerate}
  \item Eliminando de forma definitiva a necessidade de escrever casos de teste unitário, já que o processo de compilação cruzada garante a integridade incondicional das regras de negócio.

  \item Restringindo rigorosamente a arquitetura dos projetos a um limite máximo de linhas de código por pacote para evitar que os softwares atinjam uma escala insustentável.

  \item Forçando a adoção de metodologias de desenvolvimento ágil (\emph{Agile/Scrum}) através de ferramentas nativas embutidas diretamente nos compiladores das linguagens funcionais.

  \item Reduzindo o tempo de desenvolvimento ao automatizar a geração de código através de modelos probabilísticos baseados em inteligência artificial no sistema de tipos dinâmicos.

  \item Elevando o nível de abstração focado na declaração (o que fazer) e adotando bases que eliminam o estado mutável não documentado, o que reduz drasticamente a complexidade do rastreamento de defeitos.


\end{enumerate}

\questionTitle{34}{Entrada não é uma função pura}
\begin{flushright}
  \footnotesize
  \textit{\textbf{8: Programas interativos}}\\
  \textit{c08-modelo-io-04}\\

\end{flushright}
Por que uma função que lê uma linha do teclado não pode ter simplesmente o tipo \haskellinline|String|?

\begin{enumerate}
  \item Porque toda entrada digitada pelo usuário é sempre um único caractere, então o tipo correto seria \haskellinline|Char|.

  \item Porque \haskellinline|String| só pode representar literais escritos diretamente no código-fonte.

  \item Porque o compilador converte automaticamente toda leitura do teclado para \haskellinline|Int|.

  \item Porque \haskellinline|String| é um tipo numérico e não pode ser usado para texto vindo do terminal.

  \item Porque o valor lido depende de uma interação externa. Haskell representa essa interação como uma ação do tipo \haskellinline|IO String|, não como uma \haskellinline|String| pura.


\end{enumerate}

\questionTitle{35}{Desvantagem do Integer}
\begin{flushright}
  \footnotesize
  \textit{\textbf{4: Tipos de Dados}}\\
  \textit{c04-tipos-int-integer-03}\\

\end{flushright}
Qual é a principal desvantagem de usar o tipo \haskellinline{Integer} em vez de \haskellinline{Int} para todos os cálculos inteiros?

\begin{enumerate}
  \item A falta de precisão para números muito grandes.

  \item A possibilidade de ocorrer um \emph{overflow}.

  \item A incompatibilidade com operadores aritméticos básicos como \haskellinline{+}.

  \item O desempenho, pois os cálculos com \haskellinline{Integer} são geralmente mais lentos.

  \item A dificuldade de converter para outros tipos numéricos.


\end{enumerate}

\questionTitle{36}{Aninhamento com três casos}
\begin{flushright}
  \footnotesize
  \textit{\textbf{5: Expressão condicional}}\\
  \textit{c05-if-aninhamento-02}\\

\end{flushright}
Avalie a função abaixo na entrada \haskellinline|f 5|:
\begin{haskellcode}
f x = if x > 0
      then if x > 10
           then 1
           else 2
      else 3
\end{haskellcode}

\begin{enumerate}
  \item \haskellinline|2|

  \item Ocorre um erro de sintaxe por falta de parênteses.

  \item \haskellinline|3|

  \item Ocorre um erro de ambiguidade entre os dois \haskellinline|else|.

  \item \haskellinline|1|


\end{enumerate}

\questionTitle{37}{Lookup com valor padrão}
\begin{flushright}
  \footnotesize
  \textit{\textbf{6: Estruturas de dados básicas}}\\
  \textit{c06-lookup-associacao-02}\\

\end{flushright}
Qual é o resultado da expressão abaixo?
\begin{minted}{haskell}
let tabela = [("joao",3.0),("carla",9.0)]
in fromMaybe 0.0 (lookup "maria" tabela)
\end{minted}

\begin{enumerate}
  \item \haskellinline{ 0.0 }

  \item Ocorre um erro em tempo de execução.

  \item \haskellinline{ Nothing }

  \item \haskellinline{ 9.0 }

  \item \haskellinline{ Just 0.0 }


\end{enumerate}

\questionTitle{38}{Modelagem de inventário}
\begin{flushright}
  \footnotesize
  \textit{\textbf{6: Estruturas de dados básicas}}\\
  \textit{c06-modelagem-02}\\

\end{flushright}
Você quer representar um inventário como uma lista de produtos. Cada produto possui nome (\haskellinline{String}), código (\haskellinline{Int}) e preço opcional (\haskellinline{Maybe Double}). Qual tipo representa corretamente esse inventário?

\begin{enumerate}
  \item \haskellinline{[(String, Int, Maybe Double)]}

  \item \haskellinline{([String], [Int], [Maybe Double])}

  \item \haskellinline{(String, Int, Maybe Double)}

  \item \haskellinline{[String, Int, Maybe Double]}

  \item \haskellinline{Maybe (String, Int, Double)}


\end{enumerate}

\questionTitle{39}{Tipo da condição}
\begin{flushright}
  \footnotesize
  \textit{\textbf{5: Expressão condicional}}\\
  \textit{c05-if-regras-02}\\

\end{flushright}
Qual das expressões abaixo falha especificamente porque a condição do \haskellinline|if| não tem tipo \haskellinline|Bool|?

\begin{enumerate}
  \item \haskellinline|if 3 < 5 then 10 else 20|

  \item \haskellinline|if 3 then 10 else 20|

  \item \haskellinline|if even 4 then 10 else 20|

  \item \haskellinline|if null [] then 10 else 20|

  \item \haskellinline|if True then 10 else 20|


\end{enumerate}

\questionTitle{40}{A analogia da receita de bolo}
\begin{flushright}
  \footnotesize
  \textit{\textbf{1: Paradigmas de programação}}\\
  \textit{c01-conceitos-paradigmas-03}\\

\end{flushright}
Ao explicar paradigmas de programação, é comum utilizar a analogia de uma \emph{receita de bolo} (ex: \emph{1. quebre os ovos; 2. adicione farinha; 3. asse por 40 minutos}). 

Esta analogia ilustra de forma primária o raciocínio de qual paradigma e por quê?

\begin{enumerate}
  \item O paradigma orientado a objetos, pois a receita representa a troca de mensagens assíncronas entre as instâncias dos ingredientes.

  \item O paradigma declarativo, pois descreve as propriedades finais do bolo sem se preocupar com a ordem matemática em que os ingredientes foram misturados.

  \item O paradigma imperativo, pois foca na descrição passo a passo da sequência de controle e nas mudanças sucessivas de estado para atingir o objetivo.

  \item O paradigma funcional puro, pois cada passo da receita é uma função referencialmente transparente que não altera o estado da cozinha.

  \item O paradigma lógico, pois cada passo da receita atua como um axioma em um motor de inferência que deduz a existência do bolo.


\end{enumerate}

\questionTitle{41}{Conversão de operador infixo para função prefixa}
\begin{flushright}
  \footnotesize
  \textit{\textbf{3: Expressões e definições}}\\
  \textit{c03-aplicacao-notacoes-03}\\

\end{flushright}
Qual alternativa é uma reescrita prefixa equivalente à expressão \haskellinline|100 `mod` (5 * 2)|?

\begin{enumerate}
  \item \haskellinline|100 mod (5 * 2)|.

  \item \haskellinline|(mod 100 5) * 2|.

  \item \haskellinline|mod 100 (5 * 2)|.

  \item \haskellinline|mod (100, 5 * 2)|.

  \item \haskellinline|mod (100 * 5) 2|.


\end{enumerate}

\questionTitle{42}{Associatividade de subtração e potenciação}
\begin{flushright}
  \footnotesize
  \textit{\textbf{3: Expressões e definições}}\\
  \textit{c03-precedencia-associatividade-03}\\

\end{flushright}
Sabendo que \haskellinline|-| associa à esquerda e \haskellinline|^| associa à direita, qual alternativa está correta?

\begin{enumerate}
  \item Tanto \haskellinline|-| quanto \haskellinline|^| associam sempre à direita.

  \item Tanto \haskellinline|-| quanto \haskellinline|^| associam sempre à esquerda.

  \item A associatividade só se aplica a funções prefixas, não a operadores infixos.

  \item \haskellinline|10 - 5 - 2| é lida como \haskellinline|10 - (5 - 2)|, enquanto \haskellinline|2 ^ 3 ^ 2| é lida como \haskellinline|(2 ^ 3) ^ 2|.

  \item \haskellinline|10 - 5 - 2| é lida como \haskellinline|(10 - 5) - 2|, enquanto \haskellinline|2 ^ 3 ^ 2| é lida como \haskellinline|2 ^ (3 ^ 2)|.


\end{enumerate}

\questionTitle{43}{Fibonacci com acumuladores}
\begin{flushright}
  \footnotesize
  \textit{\textbf{9: Funções recursivas}}\\
  \textit{c09-cauda-fibonacci-01}\\

\end{flushright}
Em uma versão de Fibonacci com acumuladores, mantém-se um par \haskellinline|(a,b)| de números consecutivos da sequência. Qual atualização preserva corretamente essa interpretação?

\begin{enumerate}
  \item Atualizar \haskellinline|(a,b)| para \haskellinline|(b,a+b)|.

  \item Atualizar \haskellinline|(a,b)| para \haskellinline|(b,a*b)|.

  \item Atualizar \haskellinline|(a,b)| para \haskellinline|(a,b+1)|.

  \item Atualizar \haskellinline|(a,b)| para \haskellinline|(a-1,b-1)|.

  \item Atualizar \haskellinline|(a,b)| para \haskellinline|(a+b,b)|.


\end{enumerate}

\questionTitle{44}{putStrLn e prompts}
\begin{flushright}
  \footnotesize
  \textit{\textbf{8: Programas interativos}}\\
  \textit{c08-buffering-04}\\

\end{flushright}
Por que \haskellinline|putStrLn "Digite seu nome:"| costuma evitar o problema de \emph{prompt} invisível que pode ocorrer com \haskellinline|putStr "Digite seu nome: "|?

\begin{enumerate}
  \item Porque \haskellinline|putStrLn| converte o texto para \haskellinline|IO String| antes de imprimir.

  \item Porque \haskellinline|putStrLn| só funciona dentro de programas compilados, não no GHCi.

  \item Porque \haskellinline|putStrLn| acrescenta uma quebra de linha; em muitos ambientes interativos, essa quebra de linha faz a saída ser descarregada. Porém, ela também coloca a entrada do usuário na linha seguinte.

  \item Porque \haskellinline|putStrLn| executa automaticamente \haskellinline|getLine| depois de imprimir.

  \item Porque \haskellinline|putStrLn| não usa \emph{buffer}, enquanto \haskellinline|putStr| sempre usa.


\end{enumerate}

\questionTitle{45}{Composição com fromMaybe}
\begin{flushright}
  \footnotesize
  \textit{\textbf{6: Estruturas de dados básicas}}\\
  \textit{c06-maybe-operacoes-01}\\

\end{flushright}
Qual é o resultado da expressão abaixo?
\begin{minted}{haskell}
fromMaybe "vazio" (if False then Nothing else Just "ok")
\end{minted}

\begin{enumerate}
  \item \haskellinline{ Just "ok" }

  \item \haskellinline{ "ok" }

  \item \haskellinline{ Nothing }

  \item Ocorre um erro de tipos incompatíveis.

  \item \haskellinline{ "vazio" }


\end{enumerate}

\questionTitle{46}{Primeira guarda verdadeira vence}
\begin{flushright}
  \footnotesize
  \textit{\textbf{5: Expressão condicional}}\\
  \textit{c05-guardas-avaliacao-02}\\

\end{flushright}
Considere a função:
\begin{haskellcode}
f x
  | x >= 0    = 10
  | x > 3     = 20
  | otherwise = 30
\end{haskellcode}
Qual é o valor de \haskellinline|f 4|?

\begin{enumerate}
  \item \haskellinline|20|

  \item O compilador rejeita a definição pela sobreposição das guardas.

  \item \haskellinline|30|

  \item \haskellinline|10|

  \item Ocorre um erro, pois duas guardas são verdadeiras.


\end{enumerate}

\questionTitle{47}{Regra de nomenclatura de tipos}
\begin{flushright}
  \footnotesize
  \textit{\textbf{4: Tipos de Dados}}\\
  \textit{c04-conceito-nomes-01}\\

\end{flushright}
Qual é a regra sintática obrigatória para nomes de tipos em Haskell?

\begin{enumerate}
  \item Devem ser escritos inteiramente em maiúsculas.

  \item Devem começar com uma letra maiúscula.

  \item Devem começar com uma letra minúscula.

  \item Não podem conter caracteres especiais.

  \item Devem conter pelo menos um dígito.


\end{enumerate}

\questionTitle{48}{Ordem de avaliação das guardas}
\begin{flushright}
  \footnotesize
  \textit{\textbf{5: Expressão condicional}}\\
  \textit{c05-guardas-sintaxe-semantica-02}\\

\end{flushright}
Como as guardas de uma definição são avaliadas?

\begin{enumerate}
  \item O resultado é escolhido pela guarda cuja expressão à direita de \haskellinline|=| parece mais simples.

  \item O compilador reordena as guardas para testar primeiro as mais específicas.

  \item Todas as guardas são avaliadas, e a última verdadeira determina o resultado.

  \item As guardas são avaliadas de baixo para cima.

  \item De cima para baixo; a primeira guarda verdadeira determina o resultado.


\end{enumerate}

\questionTitle{49}{Domínio de uma função recursiva}
\begin{flushright}
  \footnotesize
  \textit{\textbf{9: Funções recursivas}}\\
  \textit{c09-fundamentos-dominio-01}\\

\end{flushright}
Ao definir uma função recursiva sobre números naturais em Haskell, por que é importante explicitar o domínio ou tratar entradas fora dele?

\begin{enumerate}
  \item Porque todo caso recursivo precisa ter exatamente dois casos base.

  \item Porque funções recursivas só podem ser usadas com o tipo \haskellinline|Natural|, nunca com \haskellinline|Integer|.

  \item Porque Haskell não permite definir funções recursivas sem testar todos os inteiros possíveis explicitamente.

  \item Porque o compilador sempre rejeita funções recursivas que não tratam números negativos.

  \item Porque a função pode receber valores fora do domínio matemático pretendido, e isso pode causar falhas de cobertura ou chamadas que nunca alcançam o caso base.


\end{enumerate}

\questionTitle{50}{Soma de quantidade fixa de valores}
\begin{flushright}
  \footnotesize
  \textit{\textbf{10: Ações de E/S recursivas}}\\
  \textit{c10-rastreamento-04}\\

\end{flushright}
A ação abaixo é chamada como \haskellinline|lerESomar 2|. Depois disso, o usuário digita exatamente os valores \haskellinline|10, -3|, um por linha.

\begin{haskellcode}
lerESomar :: Integer -> IO Integer
lerESomar n
  | n <= 0 = return 0
  | n > 0  = do
      x <- readLn
      s <- lerESomar (n - 1)
      return (x + s)
\end{haskellcode}

Qual valor é produzido pela ação?

\begin{enumerate}
  \item 9

  \item 2

  \item 0

  \item 6

  \item 7


\end{enumerate}

\questionTitle{51}{Assinatura de função com Char, Int e String}
\begin{flushright}
  \footnotesize
  \textit{\textbf{4: Tipos de Dados}}\\
  \textit{c04-funcao-assinatura-03}\\

\end{flushright}
Qual é a assinatura de tipo de uma função que recebe um \haskellinline{Char}, um \haskellinline{Int} e retorna uma \haskellinline{String}?

\begin{enumerate}
  \item \haskellinline{Char -> Int -> String}

  \item \haskellinline{Char, Int -> String}

  \item \haskellinline{[Char, Int, String]}

  \item \haskellinline{(Char, Int) -> String}

  \item \haskellinline{String -> Char -> Int}


\end{enumerate}

\questionTitle{52}{Algoritmo de Euclides}
\begin{flushright}
  \footnotesize
  \textit{\textbf{9: Funções recursivas}}\\
  \textit{c09-projeto-mdc-01}\\

\end{flushright}
Qual alternativa expressa corretamente o passo recursivo do algoritmo de Euclides para calcular \haskellinline|mdc a b| quando \haskellinline|b > 0|?

\begin{enumerate}
  \item Retornar imediatamente \haskellinline|a + b|.

  \item Chamar recursivamente \haskellinline|mdc (a + b) b|.

  \item Chamar recursivamente \haskellinline|mdc b (mod a b)|.

  \item Chamar recursivamente \haskellinline|mdc (a - b) (b - a)| sem verificar sinais.

  \item Chamar recursivamente \haskellinline|mdc a (mod b a)|.


\end{enumerate}

\questionTitle{53}{Eficiência e recomputação}
\begin{flushright}
  \footnotesize
  \textit{\textbf{9: Funções recursivas}}\\
  \textit{c09-eficiencia-fibonacci-01}\\

\end{flushright}
Ao comparar a definição ingênua de Fibonacci com uma versão usando acumuladores, qual é a principal diferença de eficiência discutida no capítulo?

\begin{enumerate}
  \item A versão com acumuladores só funciona para \haskellinline|fib 0| e \haskellinline|fib 1|.

  \item A versão ingênua recalcula subproblemas repetidamente, enquanto a versão com acumuladores carrega resultados parciais de forma mais direta.

  \item As duas versões sempre têm exatamente o mesmo número de chamadas.

  \item A versão ingênua é mais rápida porque usa duas chamadas recursivas.

  \item A versão com acumuladores deixa de usar casos base.


\end{enumerate}

\questionTitle{54}{Regra para abreviação de comandos}
\begin{flushright}
  \footnotesize
  \textit{\textbf{2: Ambiente de desenvolvimento}}\\
  \textit{c02-ghci-inspecao-03}\\

\end{flushright}
No ambiente GHCi, os comandos que começam com dois-pontos (\texttt{:}) frequentemente podem ser abreviados (por exemplo, \texttt{:r} no lugar de \texttt{:reload}). Qual é a regra técnica exigida pelo ambiente para que uma abreviação seja aceita?

\begin{enumerate}
  \item A abreviação deve ter exatamente duas letras para que o analisador do GHCi não as confunda com nomes de variáveis Haskell válidas.

  \item A abreviação só será aceita se o programador tiver configurado previamente um arquivo de manifesto indicando quais atalhos ele deseja mapear.

  \item O GHCi aceita qualquer abreviação independentemente de colisão; em caso de empate, ele sempre executa o último comando digitado no histórico da sessão.

  \item Apenas comandos que controlam o fluxo do terminal (como sair ou carregar) suportam atalhos; comandos de inspeção devem ser sempre digitados por extenso.

  \item A abreviação informada deve ser não ambígua, ou seja, deve ser um prefixo exclusivo que não coincida com a abreviação de nenhum outro comando existente.


\end{enumerate}

\questionTitle{55}{Projeto para cálculo de média}
\begin{flushright}
  \footnotesize
  \textit{\textbf{10: Ações de E/S recursivas}}\\
  \textit{c10-separacao-puro-03}\\

\end{flushright}
Para o exercício da média de uma sequência, considere o objetivo: ler valores até um sentinela negativo e imprimir a média dos valores não negativos lidos. Qual decomposição é mais alinhada ao estilo do capítulo?

\begin{enumerate}
  \item Uma única função pura \haskellinline|media :: IO Double| que chama \haskellinline|readLn| internamente sem usar \haskellinline|do|.

  \item Uma ação \haskellinline|lerLista :: IO [Double]| para coletar os valores e uma função pura ou expressão separada para calcular a média da lista, tratando a lista vazia antes da divisão.

  \item Uma função \haskellinline|sum :: IO [Double] -> Double|, pois \haskellinline|sum| opera diretamente sobre ações de E/S.

  \item Uma ação \haskellinline|main| que acumula uma string com todos os números e chama \haskellinline|read| apenas no final.

  \item Uma definição que ignora a lista vazia, pois \haskellinline|fromIntegral (length lista)| nunca pode ser zero.


\end{enumerate}

\questionTitle{56}{Segurança de tipos com sinônimos}
\begin{flushright}
  \footnotesize
  \textit{\textbf{4: Tipos de Dados}}\\
  \textit{c04-sinonimos-seguranca-01}\\

\end{flushright}
Considere as definições \haskellinline{type UsuarioID = Int} e \haskellinline{type ProdutoID = Int}. Se uma função espera um \haskellinline{UsuarioID} mas recebe um \haskellinline{ProdutoID}, o que acontecerá?

\begin{enumerate}
  \item Ocorrerá um erro em tempo de execução.

  \item O programa compilará normalmente, pois para o sistema de tipos, \haskellinline{UsuarioID}, \haskellinline{ProdutoID} e \haskellinline{Int} são o mesmo tipo.

  \item A palavra-chave \haskellinline{type} não pode ser usada com tipos numéricos.

  \item O programa falhará na compilação com um erro de incompatibilidade de tipos.

  \item O programa compilará, mas emitirá um aviso (\emph{warning}).


\end{enumerate}

\questionTitle{57}{Explorando módulos com :browse}
\begin{flushright}
  \footnotesize
  \textit{\textbf{2: Ambiente de desenvolvimento}}\\
  \textit{c02-ghci-modulos-03}\\

\end{flushright}
Após adicionar um módulo à sessão atual do GHCi, o desenvolvedor digita o comando \texttt{:browse Data.List}. Qual será o resultado visível no terminal?

\begin{enumerate}
  \item O GHCi abrirá uma nova aba no navegador web (\emph{browser}) do sistema operacional exibindo a documentação HTML oficial do módulo.

  \item O GHCi removerá o módulo do escopo ativo, fechando as portas para futuras avaliações matemáticas.

  \item O GHCi listará as assinaturas de tipo de todas as funções e definições de dados que estão exportadas e disponíveis publicamente naquele módulo.

  \item O GHCi procurará erros de compilação (\emph{warnings}) ignorados nas funções internas daquele módulo específico.

  \item O GHCi exibirá o código-fonte C gerado nos bastidores para cada função presente em \texttt{Data.List}.


\end{enumerate}

\questionTitle{58}{Strings como listas}
\begin{flushright}
  \footnotesize
  \textit{\textbf{6: Estruturas de dados básicas}}\\
  \textit{c06-listas-operacoes-02}\\

\end{flushright}
Como \haskellinline{String} é uma lista de caracteres, a expressão abaixo pode ser avaliada pelas mesmas funções de listas.
\begin{minted}{haskell}
drop 4 "Haskell"
\end{minted}
Qual é o resultado?

\begin{enumerate}
  \item \haskellinline{ "askell" }

  \item \haskellinline{ "ell" }

  \item \haskellinline{ "Hask" }

  \item \haskellinline{ "skell" }

  \item A expressão resulta em erro de tipo.


\end{enumerate}

\questionTitle{59}{Cabeça, cauda e construção}
\begin{flushright}
  \footnotesize
  \textit{\textbf{6: Estruturas de dados básicas}}\\
  \textit{c06-listas-estrutura-02}\\

\end{flushright}
Qual é o resultado da expressão abaixo?
\begin{minted}{haskell}
tail (['x'] ++ "yz")
\end{minted}

\begin{enumerate}
  \item \haskellinline{ "z" }

  \item \haskellinline{ "yz" }

  \item \haskellinline{ "x" }

  \item \haskellinline{ "xyz" }

  \item A expressão resulta em erro de execução.


\end{enumerate}

\questionTitle{60}{Operações pendentes na pilha}
\begin{flushright}
  \footnotesize
  \textit{\textbf{9: Funções recursivas}}\\
  \textit{c09-rastreamento-pilha-01}\\

\end{flushright}
Na definição
\begin{haskellcode}
fatorial n
  | n == 0 = 1
  | n > 0  = n * fatorial (n - 1)
\end{haskellcode}
por que a chamada \haskellinline|fatorial 4| deixa operações pendentes?

\begin{enumerate}
  \item Porque o caso base de \haskellinline|fatorial| nunca é alcançado para entradas positivas.

  \item Porque cada chamada ainda precisa multiplicar o resultado da chamada recursiva pelo valor atual de \haskellinline|n|.

  \item Porque toda função com guardas é automaticamente não terminante.

  \item Porque a multiplicação é executada antes da chamada recursiva terminar.

  \item Porque \haskellinline|Integer| força o uso de listas intermediárias.


\end{enumerate}

\questionTitle{61}{Escopo limitado à equação}
\begin{flushright}
  \footnotesize
  \textit{\textbf{5: Expressão condicional}}\\
  \textit{c05-where-escopo-02}\\

\end{flushright}
Considere a definição:
\begin{haskellcode}
f x
  | y > 0     = y
  | otherwise = 0
  where y = x - 1

f x = y + 10
\end{haskellcode}
Qual afirmação está correta?

\begin{enumerate}
  \item A segunda equação não pode usar \haskellinline|y|, pois esse nome vale apenas na primeira equação.

  \item A segunda equação pode usar \haskellinline|y| normalmente, pois \haskellinline|where| tem escopo em toda a função.

  \item Ambas as equações compartilham \haskellinline|y|, mas apenas se houver uma única assinatura de tipo.

  \item \haskellinline|where| sempre define variáveis globais no módulo.

  \item A segunda equação pode usar \haskellinline|y|, mas somente se ela também tiver guardas.


\end{enumerate}

\questionTitle{62}{Identificação de violação de pureza (I/O temporal)}
\begin{flushright}
  \footnotesize
  \textit{\textbf{1: Paradigmas de programação}}\\
  \textit{c01-transparencia-referencial-02}\\

\end{flushright}
Considere uma função \haskellinline|obterHoraAtual| que consulta o relógio interno do sistema operacional e retorna uma representação do momento exato em que foi chamada.

Por que a existência desta função viola a propriedade de pureza em um modelo matemático estrito?

\begin{enumerate}
  \item Porque seu valor de retorno depende de um estado global em mutação constante fora do controle do programa, fazendo com que chamadas com os mesmos argumentos produzam resultados diferentes.

  \item Porque o tempo é uma grandeza contínua e linguagens funcionais só conseguem modelar estruturas de dados e matrizes estritamente discretas.

  \item Porque a leitura do relógio do sistema exige chamadas de baixo nível em \emph{Assembly}, que quebram as abstrações do compilador Haskell.

  \item Ela não viola a pureza; contanto que a função não modifique a hora do sistema (apenas efetue uma leitura), ela continua sendo referencialmente transparente.

  \item Porque ela consome excessivos ciclos de processamento do \emph{garbage collector}, atrasando a execução do fluxo principal da aplicação.


\end{enumerate}

\questionTitle{63}{Let como subexpressão}
\begin{flushright}
  \footnotesize
  \textit{\textbf{3: Expressões e definições}}\\
  \textit{c03-let-escopo-03}\\

\end{flushright}
Qual é o valor da expressão \haskellinline|2 * (let x = 3; y = 4 in x + y)|?

\begin{enumerate}
  \item \haskellinline|7|.

  \item \haskellinline|24|.

  \item \haskellinline|9|.

  \item \haskellinline|14|.

  \item A expressão é inválida, pois \haskellinline|let| não pode aparecer dentro de parênteses.


\end{enumerate}

\questionTitle{64}{If como subexpressão}
\begin{flushright}
  \footnotesize
  \textit{\textbf{5: Expressão condicional}}\\
  \textit{c05-if-avaliacao-02}\\

\end{flushright}
Considere a expressão:
\begin{haskellcode}
4 * (if 'B' < 'A' then 2 + 3 else 2) - 3
\end{haskellcode}
Qual é o seu valor?

\begin{enumerate}
  \item \haskellinline|5|

  \item \haskellinline|20|

  \item Ocorre um erro de tipos.

  \item \haskellinline|17|

  \item \haskellinline|-1|


\end{enumerate}

\questionTitle{65}{Ações que exibem valores retornam IO unit}
\begin{flushright}
  \footnotesize
  \textit{\textbf{10: Ações de E/S recursivas}}\\
  \textit{c10-lacos-io-04}\\

\end{flushright}
Considere a ação:

\begin{haskellcode}
mostraLista :: Show a => [a] -> IO ()
mostraLista lista
  | null lista = return ()
  | otherwise = do
      print (head lista)
      mostraLista (tail lista)
\end{haskellcode}

Por que o tipo de retorno é \haskellinline|IO ()|?

\begin{enumerate}
  \item Porque \haskellinline|return ()| interrompe o programa e impede que qualquer valor seja produzido.

  \item Porque toda ação recursiva em Haskell deve obrigatoriamente ter tipo \haskellinline|IO ()|.

  \item Porque a lista é convertida em uma tupla vazia antes de ser percorrida.

  \item Porque \haskellinline|print| apaga os elementos da lista e sempre retorna a lista vazia.

  \item Porque a ação realiza efeitos de saída, mas não produz um valor de domínio relevante ao final.


\end{enumerate}

\questionTitle{66}{Avaliação lazy e argumento não usado}
\begin{flushright}
  \footnotesize
  \textit{\textbf{3: Expressões e definições}}\\
  \textit{c03-layout-avaliacao-03}\\

\end{flushright}
Considere a expressão:

\begin{haskellcode}
(\x -> 1) (div 1 0)
\end{haskellcode}

O que acontece em Haskell, considerando a avaliação \emph{lazy}?

\begin{enumerate}
  \item A expressão sempre produz erro de divisão por zero antes de a função ser aplicada, como ocorreria em avaliação estrita.

  \item A expressão é sintaticamente inválida, pois uma lambda não pode ser aplicada diretamente a um argumento.

  \item A expressão avalia para a própria função \haskellinline|\x -> 1|, sem aplicar o argumento.

  \item A expressão avalia para \haskellinline|1|, pois o argumento \haskellinline|div 1 0| não é necessário para produzir o resultado da função.

  \item A expressão avalia para \haskellinline|0|, pois Haskell define \haskellinline|div 1 0| como zero em contextos não usados.


\end{enumerate}

\questionTitle{67}{Homogeneidade de listas}
\begin{flushright}
  \footnotesize
  \textit{\textbf{4: Tipos de Dados}}\\
  \textit{c04-checagem-listas-01}\\

\end{flushright}
Em Haskell, se um programador tentar avaliar a expressão \haskellinline{[1, True, "texto"]}, o que acontecerá?

\begin{enumerate}
  \item A expressão será aceita e o tipo inferido será \haskellinline{Any}.

  \item O compilador criará uma tupla automaticamente.

  \item Os valores serão implicitamente convertidos para \haskellinline{String}.

  \item Ocorrerá um erro em tempo de compilação, pois todos os elementos de uma lista devem ter o mesmo tipo.

  \item Ocorrerá um erro apenas em tempo de execução, ao tentar usar os elementos.


\end{enumerate}

\questionTitle{68}{Classe mínima para div e mod}
\begin{flushright}
  \footnotesize
  \textit{\textbf{7: Polimorfismo}}\\
  \textit{cap07-hierarquia-classe-minima-divisao}\\

\end{flushright}
Considere uma função que usa apenas \haskellinline{div} e \haskellinline{mod} sobre seus argumentos. Qual restrição é a mais adequada para a assinatura mais geral?
\begin{enumerate}
  \item \haskellinline{Integral a}, pois \haskellinline{div} e \haskellinline{mod} operam sobre tipos integrais.
  \item Nenhuma restrição é necessária, pois \haskellinline{div} é paramétrico puro.
  \item \haskellinline{Floating a}, pois qualquer divisão exige ponto flutuante.
  \item \haskellinline{Fractional a}, pois divisão sempre exige tipos fracionários.
  \item \haskellinline{Show a}, pois \haskellinline{div} converte os operandos para texto.

\end{enumerate}

\questionTitle{69}{Lista final produzida por acumulador com reverse}
\begin{flushright}
  \footnotesize
  \textit{\textbf{10: Ações de E/S recursivas}}\\
  \textit{c10-listas-ordem-03}\\

\end{flushright}
A função abaixo é executada inicialmente como \haskellinline|lerLista []|. O usuário digita \haskellinline|11, 12, 0|, um valor por linha.

\begin{haskellcode}
lerLista :: [Integer] -> IO [Integer]
lerLista xs = do
  x <- readLn
  if x == 0
    then return (reverse xs)
    else lerLista (x:xs)
\end{haskellcode}

Qual lista é produzida pela ação?

\begin{enumerate}
  \item \haskellinline|[0]|

  \item \haskellinline|[]|

  \item \haskellinline|[12,11]|

  \item \haskellinline|[11,12]|

  \item \haskellinline|[11,12,0]|


\end{enumerate}


\end{multicols}
\makeanswersheet{UFOP}{BCC222: Programação Funcional}{José Romildo}{2026/1}{Primeira Prova}{05}{69}{25}{7}
\gabarito{05}{D,D,E,D,A,B,C,C,D,C,B,E,D,C,C,B,D,A,A,E,B,D,E,E,E,B,A,B,A,B,E,E,E,E,D,A,A,A,B,C,C,E,A,C,B,D,B,E,E,E,A,C,B,E,B,B,C,B,B,B,A,A,D,A,E,D,D,A,D}
\end{document}
